\chapter*{Tóm tắt}

Các bệnh lý và chấn thương xương ngày càng phổ biến, làm gia tăng nhu cầu ứng dụng trí tuệ nhân tạo 
nhằm hỗ trợ bác sĩ trong quá trình phân tích ảnh X-quang. Tuy nhiên, dữ liệu X-quang y khoa 
thường có số lượng mẫu gán nhãn hạn chế, mất cân bằng giữa các lớp bệnh và tồn tại nhiều trường hợp ngoài phân phối 
(Out-of-Distribution - OOD), khiến hiệu quả của các mô hình học sâu suy giảm khi triển khai trong thực tế.

Đề tài hướng đến xây dựng một hệ thống phân loại đa nhãn ảnh X-quang xương theo hướng đa phương thức, 
khai thác đồng thời thông tin từ ảnh X-quang và văn bản lâm sàng nhằm nâng cao hiệu quả dự đoán.

Hệ thống được xây dựng trên nền tảng BioMedCLIP với hai giai đoạn huấn luyện. 
Giai đoạn thứ nhất thực hiện thích nghi miền bằng kỹ thuật LoRA sử dụng 
Semantic Matching Loss nhằm cải thiện khả năng biểu diễn đặc trưng đối với dữ liệu X-quang xương. 
Giai đoạn thứ hai sử dụng cơ chế chú ý hai chiều để hợp nhất đặc trưng ảnh và văn bản, 
kết hợp Asymmetric Loss nhằm xử lý mất cân bằng dữ liệu và Prototypical Network để tăng cường khả năng phân biệt giữa các lớp, 
hỗ trợ các lớp có ít mẫu và làm cơ sở cho phát hiện dữ liệu ngoài phân phối. 
Hệ thống được đánh giá thông qua các thực nghiệm thành phần, 
so sánh với các mô hình nền tảng và đánh giá trên bộ dữ liệu X-quang xương được xây dựng từ dữ liệu thực tế tại 
Bệnh viện Chấn thương Chỉnh hình Thành phố Hồ Chí Minh.

Kết quả thực nghiệm cho thấy phương pháp đề xuất cải thiện hiệu quả phân loại đa nhãn so với mô hình nền, 
đồng thời nâng cao khả năng tổng quát và hỗ trợ phát hiện dữ liệu ngoài phân phối. 
Ngoài ra, đề tài còn đóng góp một bộ dữ liệu X-quang xương đã được tinh lọc và chuẩn hóa, 
tạo nền tảng cho các nghiên cứu tiếp theo về học đa phương thức và ứng dụng mô hình thị giác–ngôn ngữ trong phân tích ảnh y khoa.